\section{Conjugate factorization and an exact polynomial parity criterion}
\label{qspec:sec:anomaly}

Although $\log(x^2+b^2)=\log(x+ib)+\log(x-ib)$ for $x>0$, one cannot
split arbitrary zeta-regularized weighted sums of these logarithms without
checking the continuation. This section computes the correction rather
than assuming it vanishes.

For $a,b>0$ and a polynomial $W(x)=\sum_{d=0}^D w_dx^d\in\C[x]$, define
\begin{align}
Z_W(s;a,b)&=\sum_{n\ge0}W(n+a)Q_b(n+a)^{-s},\label{qspec:eq:weighted-quadratic}\\
F_{W,\pm}(s;a,b)&=\sum_{n\ge0}W(n+a)(n+a\pm ib)^{-s}.
\label{qspec:eq:weighted-linear}
\end{align}
Their initial absolute-convergence conditions are respectively
$\Rea s>(D+1)/2$ and $\Rea s>D+1$. We always mean their meromorphic
continuations when evaluating at zero. The sign convention for the
factorization defect is
\begin{equation}\label{qspec:eq:delta-definition}
\delta_W(a,b)=F_{W,+}'(0;a,b)+F_{W,-}'(0;a,b)-Z_W'(0;a,b).
\end{equation}
It is a convention for spectral derivatives; changing to logarithms of
regularized determinants changes the associated sign convention.

\begin{theorem}[Harmonic residue formula and universal vanishing criterion]
\label{qspec:thm:anomaly}
All three continued functions in
\eqref{qspec:eq:weighted-quadratic}--\eqref{qspec:eq:weighted-linear} are
holomorphic at zero. The defect is independent of $a$ and equals
\begin{align}
\delta_W(a,b)
&=\frac14[x^{-1}]W(x)
\left(\log(1+ib/x)-\log(1-ib/x)\right)^2\label{qspec:eq:anomaly-residue}\\
&=\sum_{k\ge1}w_{2k-1}\frac{(-1)^k b^{2k}}{k}
\left(H_{2k}-\frac12H_k\right).
\label{qspec:eq:anomaly-harmonic}
\end{align}
The last sum is finite. For a fixed polynomial $W$, the identity
$\delta_W(a,b)=0$ for every $b>0$ holds if and only if $W$ is even.
\end{theorem}
\begin{proof}
We first take $b<a$. Put $\ell_\pm(x)=\log(1\pm ib/x)$ and
$L=\ell_++\ell_-$. In a common initial half-plane, factor out $x^{-2s}$
and compare $Z_W(s)$ with
$\tfrac12(F_{W,+}(2s)+F_{W,-}(2s))$. The remaining difference is
\begin{equation}\label{qspec:eq:local-defect}
e^{-sL}-\tfrac12\bigl(e^{-2s\ell_+}+e^{-2s\ell_-}\bigr)
=-\frac{s^2}{2}(\ell_+-\ell_-)^2+O(s^3).
\end{equation}
Both the constant and first powers of $s$ cancel identically as formal
power series in $x^{-1}$.

To justify the continuation and the coefficient comparison, expand each
factor in powers of $x^{-1}$ and sum $x=n+a$. A monomial $w_dx^d$
paired with the coefficient of $x^{-j}$ contributes
$w_d\zeta(2s+j-d,a)$. For $s$ near zero, only $j=d+1$ has a pole,
and it is
\begin{equation}\label{qspec:eq:hurwitz-local-pole}
\zeta(1+2s,a)=\frac1{2s}+\gamma_0(a)+O(s).
\end{equation}
Every original positive-degree coefficient in the exponential expansion
vanishes at least to order one at $s=0$, so this possible pole is
cancelled. The degree-zero term is a finite combination of $\zeta(-d,a)$,
which is regular. The remaining infinite coefficient sum converges
normally near zero because $b/a<1$. This proves holomorphy of the
original functions at zero, not just holomorphy of their difference.

In the difference \eqref{qspec:eq:local-defect}, all regular Hurwitz
terms start at $s^2$ and have zero first derivative at zero. Multiplying
the pole in \eqref{qspec:eq:hurwitz-local-pole} by the quadratic term
in \eqref{qspec:eq:local-defect} gives
\[
Z_W'(0)-F_{W,+}'(0)-F_{W,-}'(0)
=-\tfrac14[x^{-1}]W(x)(\ell_+-\ell_-)^2.
\]
The finite part $\gamma_0(a)$ contributes only from order $s^2$ onward;
this explains both the sign and the independence from $a$.

For arbitrary $b>0$, remove a finite head with $M+a>b$ before expanding.
At each of the finitely many head terms, the first derivatives cancel
because the ordinary logarithms add with the stated branches. The tail
calculation is identical, and the residue in
\eqref{qspec:eq:hurwitz-local-pole} is still one in its spectral variable.
This proves \eqref{qspec:eq:anomaly-residue} globally.

For the coefficient formula, write
\[
\ell_+-\ell_-=2\sum_{r\ge0}\frac{(ib)^{2r+1}}{(2r+1)x^{2r+1}}.
\]
The coefficient of $x^{-2k}$ in one quarter of its square is
\[
(-1)^kb^{2k}\sum_{r=0}^{k-1}
\frac1{(2r+1)(2k-2r-1)}
=\frac{(-1)^kb^{2k}}k\sum_{r=0}^{k-1}\frac1{2r+1}.
\]
The last sum is $H_{2k}-H_k/2$. It is nonzero for every $k\ge1$.
Only odd monomials $x^{2k-1}$ can contribute to $[x^{-1}]$.
Finally the distinct powers $b^{2k}$ are linearly independent as
polynomials in $b$, proving the universal vanishing criterion.
\end{proof}

\begin{remark}[A fixed mass is different from all masses]
At a single prescribed value of $b$, different odd terms can cancel.
For example $W(x)=x^3+\tfrac23b_0^2x$ has zero defect at $b=b_0$
but is not even. Thus the theorem characterizes vanishing identically
in the mass, not an erroneous pointwise equivalence.
\end{remark}

The first defect values, with the convention
\eqref{qspec:eq:delta-definition}, are
\begin{equation}\label{qspec:eq:anomaly-examples}
\delta_1=0,\qquad \delta_x=-b^2,\qquad
\delta_{x^2}=0,\qquad \delta_{x^3}=\frac23b^4,\qquad
\delta_{x^5}=-\frac{23}{45}b^6.
\end{equation}
They are useful negative controls: deleting the anomaly from an odd-weight
calculation gives an explicitly false formula, not merely an uncertified
one.

\subsection{Finite corrected formulas for any polynomial weight}
Expand $W$ about either conjugate root by writing
\begin{equation}\label{qspec:eq:shifted-weight-coeffs}
W(y\mp ib)=\sum_{r=0}^D c_r^\pm(b)y^r.
\end{equation}
The linear-factor series are finite sums of shifted Hurwitz functions,
so Theorem~\ref{qspec:thm:anomaly} immediately gives
\begin{equation}\label{qspec:eq:general-weighted-hurwitz}
Z_W'(0;a,b)=
\sum_{r=0}^D\left(c_r^+(b)\zeta'(-r,a+ib)
+c_r^-(b)\zeta'(-r,a-ib)\right)-\delta_W(a,b).
\end{equation}
For real coefficients, the two finite sums are conjugates. Thus every
polynomially weighted quadratic logarithmic sum in this class has a
finite Hurwitz-derivative evaluation, with the exact harmonic correction
already displayed. This is a concrete specialization of the classical
local-anomaly principle, not an assumption of multiplicativity.

\subsection{Anomaly-free factorization at every negative integer}
Let
\begin{equation}\label{qspec:eq:one-sided}
Z_a(s;b)=\sum_{n\ge0}Q_b(n+a)^{-s}
\end{equation}
initially for $\Rea s>1/2$. The weight $W(x)=Q_b(x)^m$ is even for every
integer $m\ge0$, so the defect vanishes.

\begin{corollary}\label{qspec:cor:negative-hurwitz}
For $a,b>0$ and $m\ge0$,
\begin{equation}\label{qspec:eq:negative-hurwitz}
Z_a'(-m;b)=2\Rea\sum_{j=0}^m
\binom mj(-2ib)^{m-j}\zeta'(-m-j,a+ib).
\end{equation}
\end{corollary}
\begin{proof}
For this weight $Z_W(s)=Z_a(s-m;b)$. With $y=x+ib$,
\[
Q_b(x)^m=y^m(y-2ib)^m
=\sum_{j=0}^m\binom mj(-2ib)^{m-j}y^{m+j}.
\]
Insert this finite expansion into
\eqref{qspec:eq:general-weighted-hurwitz}; the defect is zero and the
second conjugate factor gives the real-part formula.
\end{proof}
No claim is made that an arbitrary higher spectral derivative factorizes
in this way. The proof uses the first-derivative cancellation and the
special even polynomial weight. At higher spectral orders, logarithmic
cross terms require their own calculation.
